% ---------------------------------------------------------------------------
% Shared visual language for every figure in this paper.
% One palette, one glyph set, one rail geometry. Figures reuse these so a
% reader who learns the notation once in Figure 1 can read all of them.
% ---------------------------------------------------------------------------
\usepackage{tikz}
\usetikzlibrary{calc,positioning,arrows.meta,decorations.pathreplacing,
                decorations.pathmorphing,fit,backgrounds,shapes.geometric,
                patterns,patterns.meta}
\usepackage{xcolor}
\usepackage{booktabs}
% newtxmath (loaded by acmart when the ACM fonts are present) already supplies
% the AMS symbols we use, and loading amssymb on top of it redefines \Bbbk.
\providecommand{\checkmark}{\ensuremath{\surd}}

% Long command strings and qualified symbol names are unbreakable words in a
% narrow column, so let TeX stretch a line rather than run it into the gutter.
\setlength{\emergencystretch}{0.8em}

% Every command line in this paper is typed literally, so the typewriter font
% must not turn -- into a dash or `' into curly quotes.
\DisableLigatures{encoding=*, family=tt*}
\usepackage{upquote}

% --- palette ---------------------------------------------------------------
\definecolor{sgtInk}{HTML}{1B1F24}
\definecolor{sgtGrey}{HTML}{8A8F96}
\definecolor{sgtWash}{HTML}{EFECE6}
\definecolor{sgtPaper}{HTML}{FBFAF8}
\definecolor{fRate}{HTML}{C0701A}   % rate limiting
\definecolor{fCache}{HTML}{14736C}  % caching
\definecolor{fRetry}{HTML}{6A4A9E}  % retry policy
\definecolor{fExport}{HTML}{A83252} % export
\definecolor{fIdx}{HTML}{3A6EA5}    % database index (colleague's work)
% Each request colour has a wash used for fills behind text, so a coloured
% region and a coloured mark can sit next to each other without competing.
\definecolor{wRate}{HTML}{FAEEE1}
\definecolor{wCache}{HTML}{E3F0EE}
\definecolor{wRetry}{HTML}{ECE7F5}
\definecolor{wExport}{HTML}{F8E7EB}
\definecolor{wIdx}{HTML}{E6EDF5}
\definecolor{wGrey}{HTML}{EEEFF0}

% --- type ------------------------------------------------------------------
\newcommand{\cmd}[1]{\texttt{\small #1}}
\newcommand{\placeholder}[1]{{\color{red}\sffamily\small [#1]}}
% the typewriter fonts have no em dash, so quoted terminal output borrows one
\newcommand{\ttdash}{\textrm{\,\textemdash\,}}
% A qualified name like api.py::handle is one unbreakable word wide enough to
% overrun a column, so allow a line to end after the file and resume at the
% function. Nothing else in the name may break.
\ExplSyntaxOn
\NewDocumentCommand \sym { m }
  { \group_begin:
    \tl_set:Nn \l_tmpa_tl {#1}
    \tl_replace_all:Nnn \l_tmpa_tl { :: } { ::\allowbreak }
    \texttt { \small \l_tmpa_tl }
    \group_end: }
\ExplSyntaxOff
\newcommand{\flab}[1]{{\sffamily\bfseries\scriptsize #1}}
\newcommand{\ftiny}[1]{{\sffamily\scriptsize #1}}
\newcommand{\fmini}[1]{{\sffamily\fontsize{5.4}{6}\selectfont #1}}
% a symbol name inside a figure, sized to sit beside 5 to 6pt sans labels
\newcommand{\fsym}[1]{{\ttfamily\fontsize{5.2}{6}\selectfont #1}}

% --- quoted terminal output ------------------------------------------------
% Every block in the walkthrough is what the command printed. The typewriter
% font has none of the marks sgt uses, so each one is drawn or borrowed from the
% roman font and the rest of the line stays in the terminal's font.
% sgt quotes shell words with a backtick and an apostrophe, and both have to
% print as the terminal prints them rather than as a matched curly pair.
\begingroup
\catcode`\'=\active \catcode`\`=\active
\gdef\sgtupquotes{\catcode39=\active \catcode96=\active
  \def'{\textquotesingle}\def`{\textasciigrave}}
\endgroup
\newenvironment{out}
  {\par\addvspace{1.0ex}\begingroup
   \ttfamily\fontsize{6.9}{8.5}\selectfont\raggedright
   \setlength{\parindent}{0pt}%
   \sgtupquotes
   \begin{list}{}{\setlength{\leftmargin}{2.0mm}\setlength{\rightmargin}{0pt}%
                  \setlength{\itemindent}{0pt}\setlength{\labelsep}{0pt}%
                  \setlength{\labelwidth}{0pt}\setlength{\topsep}{0pt}%
                  \setlength{\partopsep}{0pt}\setlength{\parsep}{0pt}%
                  \setlength{\itemsep}{0pt}}%
   \item[]}
  {\end{list}\endgroup\addvspace{1.0ex}}
\newcommand{\ttdot}{\textrm{\,\textperiodcentered\,}}
\newcommand{\ttdisc}{\textrm{\textbullet}}
\newcommand{\ttok}{\textrm{$\checkmark$}}
\newcommand{\ttundo}{\textrm{$\curvearrowleft$}}
\newcommand{\ttnew}{\textrm{$\circ$}}
\newcommand{\ttfull}{\textcolor{sgtInk}{\rule[-0.04ex]{0.62ex}{1.0ex}}\hspace{0.32ex}}
% A column of the history in which a checkpoint has no edits at all. sgt prints a
% middle dot there rather than a space, so a bar is as wide as the history and not
% as wide as the checkpoint, and two checkpoints' bars can be read against each
% other. Set in a box of the block's own width so that they line up.
% Raised to the blocks' own centre: they run from -0.04ex to 0.96ex, and a period
% centred on the x-height sits visibly below that.
\newcommand{\ttempty}{\makebox[0.62ex][c]{%
  \raisebox{0.2ex}{\textcolor{sgtGrey!70}{\textrm{\textperiodcentered}}}}\hspace{0.32ex}}
\newcommand{\ttwarn}{\tikz[baseline=-0.3mm]{%
  \draw[line width=0.35pt, color=sgtInk] (0,0)--(1.35mm,0)--(0.675mm,1.25mm)--cycle;}}
\newcommand{\ttfork}{\textrm{$\pitchfork$}}
\newcommand{\ttleft}{\textrm{$\blacktriangleleft$}}
\newcommand{\ttright}{\textrm{$\blacktriangleright$}}
\newcommand{\ttarrow}{\textrm{$\rightarrow$}}
\newcommand{\ttrule}{\textcolor{sgtGrey}{\rule[0.45ex]{1.6ex}{0.35pt}}}
% the terminal draws its tree with box-drawing characters no font in this
% document has, so each one is a pair of rules at the same weight
\newcommand{\ttvert}{\tikz[baseline=0pt]{%
  \draw[line width=0.3pt, color=sgtGrey] (0,-1.2mm) -- (0,2.1mm);}}
\newcommand{\ttbranch}{\tikz[baseline=0pt]{%
  \draw[line width=0.3pt, color=sgtGrey] (0,-1.2mm) -- (0,2.1mm);
  \draw[line width=0.3pt, color=sgtGrey] (0,0.55mm) -- (1.7mm,0.55mm);}\hspace{0.25mm}}
\newcommand{\ttlast}{\tikz[baseline=0pt]{%
  \draw[line width=0.3pt, color=sgtGrey] (0,2.1mm) -- (0,0.55mm) -- (1.7mm,0.55mm);}%
  \hspace{0.25mm}}

% --- geometry --------------------------------------------------------------
% One op glyph is 2\opsz across, and the thumbnails at the end of this file
% redeclare all three of these at about three fifths the size, because a
% thumbnail packs its marks closer together than a full page figure does and at
% full size the neighbours touch and read as one bar rather than three edits.
\newlength{\opsz}\setlength{\opsz}{1.15mm}
\newlength{\oprad}\setlength{\oprad}{1.32mm}
\newcommand{\ckpth}{1.55}

\tikzset{
  rail/.style        = {line width=1.5pt, color=#1!22, line cap=round},
  rail/.default      = sgtGrey,
  raildim/.style     = {line width=1.5pt, color=sgtGrey!18, line cap=round},
  commitcol/.style   = {line width=0.35pt, color=sgtGrey!55, densely dotted},
  nowrule/.style     = {line width=0.6pt, color=sgtInk!70},
  dep/.style         = {line width=0.35pt, color=sgtGrey!70,
                        -{Straight Barb[length=1.1mm,width=1.1mm]}},
  % An arrowhead in these figures means a direction, so a relation that holds
  % equally both ways is drawn without one. Figure 2 pairs this with the same arc
  % dotted, and the reader then has one shape whose solid form is a relation the
  % system recorded and whose dotted form is one it did not.
  share/.style       = {line width=0.35pt, color=sgtGrey!70},
  card/.style        = {rounded corners=0.6mm, line width=0.4pt},
  note/.style        = {font=\sffamily\fontsize{5.6}{6.4}\selectfont, color=sgtInk!75,
                        inner sep=0.3mm},
  lbl/.style         = {font=\sffamily\fontsize{6.2}{7}\selectfont, color=sgtInk!85,
                        inner sep=0.3mm},
  hd/.style          = {font=\sffamily\bfseries\fontsize{6.6}{7.4}\selectfont,
                        color=sgtInk, inner sep=0.3mm},
}

% --- op glyphs -------------------------------------------------------------
% Five kinds, told apart by shape so they survive greyscale printing.
% add: filled square.  extend: filled square with an open left edge.
% rework: filled disc.  prune: hollow square with a slash.  move: right triangle.
\newcommand{\opadd}[3]{\fill[#3] ($(#1,#2)+(-\opsz,-\opsz)$) rectangle ($(#1,#2)+(\opsz,\opsz)$);}
\newcommand{\opext}[3]{%
  \fill[#3] ($(#1,#2)+(-\opsz,-\opsz)$) rectangle ($(#1,#2)+(\opsz,\opsz)$);
  \fill[sgtPaper] ($(#1,#2)+(-\opsz,-0.42mm)$) rectangle ($(#1,#2)+(-0.28mm,0.42mm)$);}
\newcommand{\oprew}[3]{\fill[#3] (#1,#2) circle (\oprad);}
\newcommand{\opprune}[3]{%
  \draw[#3, line width=0.45pt, fill=sgtPaper]
    ($(#1,#2)+(-\opsz,-\opsz)$) rectangle ($(#1,#2)+(\opsz,\opsz)$);
  \draw[#3, line width=0.45pt] ($(#1,#2)+(-0.8mm,-0.8mm)$) -- ($(#1,#2)+(0.8mm,0.8mm)$);}
\newcommand{\opmove}[3]{%
  \fill[#3] ($(#1,#2)+(-\opsz,-1.25mm)$) -- ($(#1,#2)+(-\opsz,1.25mm)$)
            -- ($(#1,#2)+(1.45mm,0)$) -- cycle;}
\newcommand{\opghost}[3]{%
  \draw[#3, draw opacity=0.55, line width=0.45pt, densely dotted, fill=sgtPaper]
    ($(#1,#2)+(-\opsz,-\opsz)$) rectangle ($(#1,#2)+(\opsz,\opsz)$);}
% a checkpoint: a tick that crosses the rail
\newcommand{\ckpt}[3]{\draw[#3, line width=0.55pt] (#1,{#2-\ckpth}) -- (#1,{#2+\ckpth});}

% --- diff hunk stripe, for drawing a git commit as the thing it is ---------
% \hunk{x}{y}{width}{color}
\newcommand{\hunk}[4]{\fill[#4!85, rounded corners=0.25mm] (#1,#2) rectangle ($(#1,#2)+(#3,0.85mm)$);}

% ---------------------------------------------------------------------------
% Shared figure devices. Four of them, reused in every figure, so a reader who
% works out what a tab or a hand means in Figure 1 never has to work it out
% again. All of them assume a picture with x=1mm, y=1mm.
% ---------------------------------------------------------------------------

% (1) A panel: a rounded card lifted off the page by a soft offset shadow, with
% a coloured tab straddling its top-left corner that names it.
% \panel{x}{y}{width}{height}{colour}{TAB TEXT}
\newcommand{\panel}[6]{%
  \fill[sgtInk!8, rounded corners=1.2mm]
    ({#1+0.55},{#2-0.55}) rectangle ({#1+#3+0.55},{#2+#4-0.55});
  \fill[sgtPaper, rounded corners=1.2mm] (#1,#2) rectangle ({#1+#3},{#2+#4});
  \draw[rounded corners=1.2mm, line width=0.3pt, color=sgtInk!15]
    (#1,#2) rectangle ({#1+#3},{#2+#4});
  \node[anchor=south west, inner xsep=1.15mm, inner ysep=0.6mm, fill=#5,
        rounded corners=0.55mm, text=white,
        font=\sffamily\bfseries\fontsize{5.3}{6}\selectfont]
    at ({#1+1.3},{#2+#4-1.25}) {#6};}

% A panel with no tab, for a region that needs the card shape without a name.
\newcommand{\plainpanel}[5]{%
  \fill[sgtInk!8, rounded corners=1.2mm]
    ({#1+0.55},{#2-0.55}) rectangle ({#1+#3+0.55},{#2+#4-0.55});
  \fill[#5, rounded corners=1.2mm] (#1,#2) rectangle ({#1+#3},{#2+#4});
  \draw[rounded corners=1.2mm, line width=0.3pt, color=sgtInk!15]
    (#1,#2) rectangle ({#1+#3},{#2+#4});}

% (2) Who wrote it. The developer writes the request and the agent writes the
% code, and that split is the paper's subject, so it gets a mark of its own.
% \whodev{x}{y}{colour}   \whoagent{x}{y}{colour}
\newcommand{\whodev}[3]{%
  \fill[#3] (#1,{#2+0.52}) circle (0.5);
  \begin{scope}
    \clip ({#1-1.05},{#2-1.1}) rectangle ({#1+1.05},{#2-0.16});
    \fill[#3] (#1,{#2-0.2}) circle (0.9);
  \end{scope}}
\newcommand{\whoagent}[3]{%
  \draw[color=#3, line width=0.26pt] (#1,{#2+1.28}) -- (#1,{#2+0.92});
  \fill[#3] (#1,{#2+1.38}) circle (0.14);
  \fill[#3, rounded corners=0.24mm]
    ({#1-0.82},{#2-0.72}) rectangle ({#1+0.82},{#2+0.92});
  \fill[sgtPaper] ({#1-0.4},{#2+0.24}) circle (0.155);
  \fill[sgtPaper] ({#1+0.4},{#2+0.24}) circle (0.155);
  \fill[sgtPaper] ({#1-0.38},{#2-0.32}) rectangle ({#1+0.38},{#2-0.2});}

% (3) An annotation: a thin curved leader from the thing to a short label.
% \annot{from x}{from y}{to x}{to y}{bend}{colour}
\newcommand{\annot}[6]{%
  \draw[color=#6!70, line width=0.32pt, bend left=#5,
        -{Straight Barb[length=0.85mm,width=0.9mm]}] (#1,#2) to (#3,#4);}
\tikzset{
  acc/.style     = {font=\sffamily\bfseries\fontsize{5.4}{6.2}\selectfont,
                    color=#1, inner sep=0.3mm},
  acc/.default   = sgtInk,
  tag/.style     = {font=\sffamily\fontsize{5.2}{6}\selectfont, inner xsep=0.7mm,
                    inner ysep=0.35mm, rounded corners=0.4mm},
}

% (4) A highlighted span inside a line of code, the way an editor shows one.
% \hispan{x}{y}{width}{wash colour}
\newcommand{\hispan}[4]{%
  \fill[#4, rounded corners=0.3mm] (#1,{#2-0.75}) rectangle ({#1+#3},{#2+1.25});}

% ---------------------------------------------------------------------------
% Operation thumbnails. Each is 12mm wide and sits on the text baseline, so
% the same picture can head a paragraph in the walkthrough and appear in the
% grid in Section 5. Every thumbnail shows the state before the command on the
% left and the state after it on the right.
% ---------------------------------------------------------------------------
% A thumbnail is 12.6mm wide and has to hold three marks on each side of an
% arrow, so the marks are drawn at three fifths the size they take on a full
% page. What survives the reduction is the silhouette, which is what tells one
% kind of edit from another, and what the reduction buys is a gap of four
% tenths of a millimetre between neighbours, so three edits look like three.
\newenvironment{thumb}%
  {\begin{tikzpicture}[baseline=-0.6mm, x=1mm, y=1mm, line cap=round]%
   \setlength{\opsz}{0.7mm}\setlength{\oprad}{0.8mm}%
   \renewcommand{\ckpth}{0.94}}%
  {\end{tikzpicture}}

\newcommand{\thumbarrow}{\draw[-{Straight Barb[length=0.8mm,width=0.9mm]},
   line width=0.3pt, color=sgtGrey] (5.6,0) -- (7.0,0);}

% save: edits sitting in the working files (dotted) become recorded ops, and
% the moment the developer looked at them is marked with a checkpoint tick.
\newcommand{\thsave}{\begin{thumb}
  \draw[rail=fCache] (0,0) -- (5.0,0);
  \opadd{1.0}{0}{fCache}\opghost{2.8}{0}{fCache}\opghost{4.6}{0}{fCache}
  \thumbarrow
  \draw[rail=fCache] (7.6,0) -- (12.9,0);
  \opadd{8.3}{0}{fCache}\opadd{10.1}{0}{fCache}\opadd{11.9}{0}{fCache}
  \ckpt{12.9}{0}{fCache}
\end{thumb}}

\newcommand{\threvert}{\begin{thumb}
  \draw[rail=fCache] (0,0) -- (5,0);
  \opadd{1.0}{0}{fCache}\opadd{2.8}{0}{fCache}\opadd{4.6}{0}{fCache}
  \thumbarrow
  \draw[rail=fCache] (7.6,0) -- (12.8,0);
  \opadd{8.4}{0}{fCache}\opghost{10.4}{0}{fCache}\opghost{12.4}{0}{fCache}
\end{thumb}}

% A dotted outline spends its gap on its own stroke, so two ghosts side by side
% need a wider pitch than two filled squares do before they stop reading as one
% box with a line down the middle.
\newcommand{\threstore}{\begin{thumb}
  \draw[rail=fCache] (0,0) -- (5,0);
  \opadd{0.6}{0}{fCache}\opghost{2.6}{0}{fCache}\opghost{4.6}{0}{fCache}
  \thumbarrow
  \draw[rail=fCache] (7.6,0) -- (12.6,0);
  \opadd{8.4}{0}{fCache}\opadd{10.2}{0}{fCache}\opadd{12.0}{0}{fCache}
\end{thumb}}

\newcommand{\thfork}{\begin{thumb}
  \draw[rail=fCache] (0,0) -- (3.4,0); \opadd{1.6}{0}{fCache}
  \draw[raildim] (3.4,0) -- (5.0,1.6); \draw[raildim] (3.4,0) -- (5.0,-1.6);
  \opadd{5.6}{1.6}{fCache} \opadd{5.6}{-1.6}{fIdx}
  \node[note] at (8.0,0) {\fmini{?}};
  \draw[rail=fCache] (9.4,0) -- (12.6,0); \oprew{11.0}{0}{fCache}
\end{thumb}}

\newcommand{\thdiff}{\begin{thumb}
  \draw[rail=fCache] (0,0) -- (5,0);
  \opadd{1.2}{0}{fCache}\opadd{3.0}{0}{fCache}
  \draw[rail=fRate] (7.6,0) -- (12.6,0);
  \opadd{8.6}{0}{fRate}\opadd{10.4}{0}{fCache}\opadd{12.2}{0}{fRate}
  \draw[line width=0.3pt, color=sgtGrey] (6.3,-1.8) -- (6.3,1.8);
\end{thumb}}

% show: an edit on a rail, and the request it was filed under, read back out
\newcommand{\thwhy}{\begin{thumb}
  \draw[rail=fCache] (0,0) -- (6.6,0);
  \opadd{1.2}{0}{fCache}\opadd{3.2}{0}{fCache}\opadd{5.2}{0}{fCache}
  \thumbarrow
  \draw[fCache!60, line width=0.35pt, fill=sgtPaper, rounded corners=0.4mm]
    (7.8,-1.7) rectangle (12.6,1.7);
  \fill[fCache!40] (8.5,0.35) rectangle (11.9,0.75);
  \fill[fCache!40] (8.5,-0.75) rectangle (11.1,-0.35);
\end{thumb}}

% intent list: one feature's rail, and the cut that divides the edits made in one
% stretch of work from the edits made in the next. The gap at the cut is twice
% the gap between two edits inside a stretch, so the two runs read as two runs
% before the reader has taken in anything else about the row.
\newcommand{\thckpt}{\begin{thumb}
  \draw[rail=fRetry] (0,0) -- (12.0,0);
  \opadd{0.9}{0}{fRetry}\opadd{2.7}{0}{fRetry}\opadd{4.5}{0}{fRetry}
  \ckpt{6.0}{0}{fRetry}
  \opadd{7.5}{0}{fRetry}\opadd{9.3}{0}{fRetry}\opadd{11.1}{0}{fRetry}
\end{thumb}}

% The destination branch is feat/export, so the rail that comes live is the export
% colour. Colour carries the request throughout this paper, and a switch row drawn
% in the caching colour would say the developer switched to the caching.
\newcommand{\thswitch}{\begin{thumb}
  \draw[rail=fRate] (0,1.5) -- (5.4,1.5); \opadd{1.6}{1.5}{fRate}\opadd{3.6}{1.5}{fRate}
  \draw[raildim] (0,-1.5) -- (5.4,-1.5); \opghost{1.6}{-1.5}{fExport}\opghost{3.6}{-1.5}{fExport}
  \thumbarrow
  \draw[raildim] (7.6,1.5) -- (12.6,1.5); \opghost{9.0}{1.5}{fRate}\opghost{11.0}{1.5}{fRate}
  \draw[rail=fExport] (7.6,-1.5) -- (12.6,-1.5); \opadd{9.0}{-1.5}{fExport}\opadd{11.0}{-1.5}{fExport}
\end{thumb}}

% propose land --subset: one of two finished features crosses the line into the
% shared branch, and the other stays behind on the developer's side of it.
% The named subset is the export, so the rail that crosses is the export colour and
% the shorter rail left behind is the half-finished retry policy, which is the pair
% Section 5.4 walks through. The rail lengths carry finished against unfinished.
\newcommand{\thsubset}{\begin{thumb}
  \draw[rail=fExport] (0,1.6) -- (7.2,1.6);
  \opadd{1.4}{1.6}{fExport}\opadd{3.4}{1.6}{fExport}
  \draw[rail=fRetry] (0,-1.6) -- (5.0,-1.6);
  \opadd{1.4}{-1.6}{fRetry}\opadd{3.4}{-1.6}{fRetry}
  \draw[nowrule] (8.4,-2.8) -- (8.4,2.8);
  \draw[rail=fExport] (8.4,1.6) -- (12.6,1.6);
  \opadd{9.8}{1.6}{fExport}\opadd{11.6}{1.6}{fExport}
\end{thumb}}

% ---------------------------------------------------------------------------
% Text macros
% ---------------------------------------------------------------------------
\newcommand{\sgt}{\textsf{sgt}}
